\chapter{Introduction}
\label{ch:xintro}

The thesis built a simulator of \sgra{} to \emph{measure} things: the size and
shape of the shadow and photon ring against the Event Horizon Telescope
\citep{EHT2022}, the mass and distance from the orbit of S2
\citep{Gillessen2017, Gravity2018, Gravity2020, Gravity2022}, and a
free-quadrupole test of the no-hair theorem. Its images were instruments ---
bisected boundaries, sub-ring decompositions, flux ratios --- and the renders
that accompanied them were illustrations.

This report asks a different question of the same spacetime: \emph{what
would one actually see?} Answering it honestly takes physics the thesis did
not need. A camera does not record bolometric intensity but a spectrum
filtered by its colour response. A moving observer sees a different sky from
a static one at the same point. The light reaching the camera now left its
source at different times along different paths, so a variable source ---
the near-infrared flares of \sgra{} --- appears smeared, delayed and echoed.
And the real accretion flow of \sgra{} is not the thin disk of the thesis'
renders at all, but a hot, optically thin plasma best seen at millimetre
wavelengths.

\section{What this extension adds}

Table~\ref{tab:phenomena} lists the phenomena the simulator now shows, the
gate of the extension plan that added each one and the chapter that
documents it. Every scene uses one spacetime --- \sgra{} with $\astar = 0.94$
\citep{EHT2022} and the GRAVITY mass and distance \citep{Gravity2022}. The
only exception, a traversable wormhole, appears in an appendix as a labelled
test of the code's generality, never as a model of \sgra{}.

\begin{table}
\centering
\caption{Phenomena added by the extension. ``Thesis'' marks what the thesis
already modelled; chapters in brackets are pending.}
\label{tab:phenomena}
\small
\begin{tabular}{p{7.4cm}cc}
\toprule
Phenomenon & Gate & Chapter \\
\midrule
Shadow, lensing, photon sub-rings, $g^4$ beaming, S2 orbit in full GR,
no-hair test & thesis & --- \\
True visual colour of a redshifted blackbody; Doppler hue shift & G13 & \ref{ch:xspectral} \\
High-dynamic-range display; a thin disk at \sgra's measured accretion rate & G14 & \ref{ch:xspectral} \\
Lensed real Milky Way panorama & G15 & \ref{ch:xspectral} \\
Horizon-penetrating Kerr--Schild integrator & G16 & \ref{ch:xks} \\
Moving observers: aberration, sky Doppler shift, time dilation, frame
dragging, light that visited the ergoregion & G17 & \ref{ch:xobs} \\
Orbiting hot spots (the GRAVITY flares) with slow light and photon-ring
echoes & G18 & \ref{ch:xflares} \\
Hot accretion flow at 230\,GHz; multi-band views & G19 & \ref{ch:xriaf} \\
Polarisation & G20 & \ref{ch:xriaf} \\
Forecasts: extended-mass bias, Lense--Thirring detectability & G22 & \ref{ch:xforecast} \\
A fall through the horizon of \sgra{} & G24 & \ref{ch:xinfall} \\
Traversable wormhole (exotic, not \sgra) & G21 & App.~\ref{app:xwormhole} \\
\bottomrule
\end{tabular}
\end{table}

\section{How the new physics is validated}

The extension keeps the thesis' rule that a picture is not evidence. Each
new capability follows the same pattern:

\begin{enumerate}
\item a \emph{CPU twin} in double precision, built on the thesis' physics
library, produces the reference numbers;
\item \emph{tests} pin the twin to closed forms wherever one exists --- the
Schwarzschild limit, Bardeen--Press--Teukolsky circular orbits
\citep{BardeenPressTeukolsky1972}, Synge's shadow \citep{Synge1966}, the
Carter constants \citep{Carter1968}, exact symmetries;
\item the \emph{live GPU path} (single precision, real time) is compared
with the twin pixel by pixel, and the agreement is quoted;
\item every number in this report is written by a tool and reaches the text
through a generated macro.
\end{enumerate}

The comparisons in step~3 are not formalities. The observer comparison of
Chapter~\ref{ch:xobs} exposed a defect in the live renderer that had shifted
the prograde edge of the high-spin shadow by $\XLegacyBudgetPx$ pixels.
While building the slow-light test of Chapter~\ref{ch:xflares}, a
reimplementation of the Carter integrator without the thesis' constraint
projection disagreed with the Kerr--Schild path at the $10^{-4}$ level, which
showed that the projection matters even at moderate inclinations.

\section{Conventions}

Units are geometric, $G = c = M = 1$, unless stated. At the GRAVITY mass and
distance one gravitational radius subtends $\UasPerM\,\muas$ and one unit of
coordinate time lasts \XMinPerM\ minutes. Cameras trace the photon they
\emph{receive} (thesis Appendix on image handedness); the default camera
looks at the hole with image-right along the local $+\hat\phi$ direction and
image-up toward the spin axis, so prograde material approaches on the left.
Inclinations are measured from the spin axis; $i > 90\degr$ views the hole
from below its equatorial plane, where prograde motion appears clockwise.

\section{Outline}

Chapter~\ref{ch:xspectral} treats colour and brightness: what a camera sees
of a redshifted blackbody, the corrected live colour pipeline, the display
chain, the hypothetical thin disk fed at \sgra's real accretion rate, and the
lensed Milky Way. Chapter~\ref{ch:xks} describes the Kerr--Schild integrator
that now drives every live image. Chapter~\ref{ch:xobs} adds observers in
motion. Chapter~\ref{ch:xflares} models the GRAVITY flares as orbiting hot
spots and compares them with the observations. Chapter~\ref{ch:xriaf}
replaces the hypothetical thin disk by \sgra's real, hot accretion flow at
230\,GHz and compares its total-intensity and polarised images with the
EHT's. The remaining chapters --- forecasts for future orbit monitoring and a
fall into \sgra{} --- will follow their gates.
